\documentclass[11pt,a4paper]{article}
\usepackage[T1]{fontenc}
\usepackage{newtxtext}
\usepackage{amsmath,amsthm,mathtools}
\usepackage{newtxmath}
\usepackage[scaled=.89]{inconsolata}
\usepackage[margin=24mm]{geometry}
\usepackage{microtype,booktabs,array,longtable,xcolor,listings,fancyhdr,hyperref,bookmark,needspace}
\definecolor{inkblue}{RGB}{28,57,77}
\definecolor{codeback}{RGB}{247,248,249}
\hypersetup{colorlinks=true,linkcolor=inkblue,urlcolor=inkblue,pdftitle={Pentagon packing: a mathematical code commentary}}
\pagestyle{fancy}\fancyhf{}\fancyhead[L]{\small\textsc{Pentagon packing}}\fancyhead[R]{\small Code v003 / manuscript v005}\fancyfoot[C]{\thepage}
\setlength{\headheight}{14pt}\setlength{\parskip}{5pt}\setlength{\parindent}{0pt}\setlength{\emergencystretch}{2em}
\setcounter{tocdepth}{1}\widowpenalty=10000\clubpenalty=10000
\lstdefinestyle{proofcode}{language=[Objective]Caml,basicstyle=\ttfamily\fontsize{8}{10}\selectfont,
keywordstyle=\color{inkblue}\bfseries,commentstyle=\color{black!55},stringstyle=\color{black!55},
backgroundcolor=\color{codeback},frame=single,rulecolor=\color{black!16},framerule=.3pt,
numbers=left,numberstyle=\tiny\ttfamily,numbersep=7pt,xleftmargin=12pt,framexleftmargin=4pt,
columns=fullflexible,keepspaces=true,showstringspaces=false,breaklines=true,breakatwhitespace=false,tabsize=2}
\lstset{style=proofcode}
\newcommand{\code}[1]{\texttt{\detokenize{#1}}}
\newcommand{\ak}{a_{\mathrm{crit}}}\newcommand{\amin}{a_{\min}}
\newcommand{\eN}{\epsilon_{\mathcal N}}\newcommand{\eM}{\epsilon_{\mathcal M}}
\newcommand{\area}{\operatorname{area}}\newcommand{\arc}{\operatorname{arc}}
\newcommand{\cc}{\mathbf c}\newcommand{\RR}{\mathbb R}\newcommand{\ZZ}{\mathbb Z}
\newtheorem{proposition}{Proposition}[section]
\theoremstyle{remark}\newtheorem*{contract}{Implementation contract}\newtheorem*{auditnote}{Audit note}
% Generated from the compiled manuscript label map.
\newcommand{\PaperPseudoOne}{5.1}
\newcommand{\PaperPseudoTwo}{5.2}
\newcommand{\PaperPseudoArea}{5.4}
\newcommand{\PaperPseudoThree}{5.5}
\newcommand{\PaperIsosceles}{5.10}
\newcommand{\PaperThree}{5.11}
\newcommand{\PaperFour}{5.12}
\newcommand{\PaperLarge}{5.14}
\newcommand{\PaperContactCoordinates}{A.2}

\begin{document}
\begin{center}
{\color{inkblue}\Large\bfseries Packings of regular pentagons}\\[5pt]
{\large A mathematical commentary on the verification code}\\[8pt]
\code{pentagon_code_v003}\qquad\code{pentagon_v005.tex}\\[5pt]
{\small 6 November 2026}
\end{center}
\begin{abstract}\noindent
This document follows the verification from regular-pentagon contacts to
finite-cluster exclusions. It explains the coordinate conventions, interval
rejection tests, automatic derivatives, and meet-in-the-middle retention
invariant. The four contact-sign calculations are derived independently from
actual pentagon vertices and placed before the squeezing area inequality.
Source listings are extracted from the accompanying files with their actual
line numbers. Printed paper numbers, stable TeX labels, and seven-letter
source comments provide complementary navigation.

The paper supplies the geometric reductions; the code exhausts the resulting
parameter boxes. Exact identities, finite function-contract tests, numerical
diagnostics, and full interval searches have different logical roles and are
identified accordingly. Execution confirmation is generated from source-matching
logs in \code{verification/CONFIRMATION.md}, rather than asserted by a fixed
paragraph in this mathematical explanation.
\end{abstract}
\tableofcontents
\clearpage

\section{The target and the division of responsibilities}

The paper proves that a packing of congruent regular pentagons has density at
most $(5-\sqrt5)/3$, with a unique periodic equality packing. The code does not
search all infinite packings. The geometric arguments reduce the theorem to
inequalities for a single triangle or a small cluster, with specified contacts,
side bounds, and attachment relations. The code excludes a closed set of
counterexamples to each such inequality.

There are four distinct assertions to keep separate. The paper supplies a
valid reduction to the terminal cases. A coordinate map covers those cases.
The interval implementation encloses its geometric functions and never rejects
a feasible input. Finally, a terminating run exhausts the encoded
counterexample boxes. A successful run is evidence for the last assertion; it
cannot repair a missing input case or an incorrect rejection rule.

The accompanying \code{docs/CALCULATIONS.md} and \code{docs/index.html}
are reader lookup tools. They use the printed numbers from the compiled
manuscript, the stable TeX labels, and the seven-letter source-navigation codes.
No code identifier needs to be printed in the published article. A reader of a
printed lemma can look up that number and reach its source functions and run
command. The ledger \code{registry.json} is also the source for job selection.

\begin{center}\small
\begin{tabular}{@{}lll@{}}\toprule
Paper & Stable label & Existing identifier\\\midrule
Lemma \PaperPseudoOne & \code{calc:pseudo1} & \code{NKQNXUN}\\
Lemma \PaperPseudoTwo & \code{calc:pseudo2} & \code{RWWHLQT}\\
Lemma \PaperPseudoArea & \code{calc:pseudo-area} & \code{BXZBPJW}\\
Lemma \PaperPseudoThree & \code{calc:pseudo-area3} & \code{JQMRXTH}\\
Lemma \PaperIsosceles & \code{calc:dimer-isosc} & \code{KUGAKIK}\\
Lemma \PaperThree & \code{calc:pent3} & \code{HUQEJAT}\\
Lemma \PaperFour & \code{calc:pent4} & \code{QPJDYDB}\\
Lemma \PaperLarge & \code{calc:large} & \code{FHBGHHY}\\\bottomrule
\end{tabular}
\end{center}

\subsection{Files and execution}
\code{pent.ml} gives constants, interval geometry, contact charts, and simple
subdivision. \code{pet.ml} supplies a necessary test for nonoverlap of a pair.
\code{banana.ml} inverts contact coordinates and reconstructs triangles.
\code{dimer.ml} supplies the two-triangle exclusions and squeezing predicates.
\code{autodiff.ml} lifts these formulas to jets. \code{meet-in-middle.ml}
implements matching and pruning; \code{mitm-calcs.ml} supplies the particular
finite-cluster domains.

The loader evaluates definitions and inexpensive range assertions,
not proof searches. Long calculations are explicit jobs rather than side effects of loading.
Each selected job runs in a separate OCaml bytecode process. This is important
for the large mutable peripheral hash tables: processes do not share them.
Concurrency is only between independent calculations. It changes neither the
internal order of a calculation nor the numerical expressions it evaluates.

\code{python3 run.py --tests} selects the deterministic contract suite.
\code{--extended} adds seeded generated diagnostics. A job can be selected by
\code{--job}; \code{--claim 5.2}, \code{--claim RWWHLQT}, and
\code{--claim calc:pseudo2} select the same direct calculation association.
\code{--section 5} selects the associated principal lemma calculations.
The listing option performs no calculation. The selector does not silently
expand the entire mathematical dependency graph of cited lemmas; use
\code{--all} for all current proof jobs, and add \code{--diagnostics} for the
optional historical and supplementary computations.

\section{Intervals, inequalities, and failure}

The normalization is
\[
\kappa=\cos(\pi/5),\qquad \sigma=\sin(\pi/5),\qquad
\ak=\frac32\kappa\sigma(1+\kappa),\qquad \amin=1.237,
\]
with $\eN=\ak-\amin$ and $\eM=0.008$. The code names are
\code{kappa}, \code{sigma}, \code{aK} or \code{acrit}, \code{amin},
\code{epsN_I}, and \code{epsM_I}. Rational constants are enclosed through
integer conversion and interval division; a machine-float singleton is not
silently described as the exact decimal rational.

An interval is a closed pair of endpoints, stored as \code{low}, \code{high}.
The relation \code{x >> y} means $\inf x>\sup y$. In particular, equality
with a target cannot satisfy a strict rejection. The definition of
\code{disjoint_I} uses strict separation in either direction, so touching
intervals are not disjoint.

\lstinputlisting[firstline=112,lastline=116,firstnumber=112,caption={\code{ocaml/pent.ml}: \code{disjoint_I}}]{../ocaml/pent.ml}

This direction of comparison is central to the proof. An area interval whose
lower endpoint exceeds the upper endpoint of the target certifies that the
entire box is outside the counterexample set. Failure to prove that inequality
means only that the box needs more work.

\subsection{The compatibility layer}
The release uses \code{Interval_crlibm.I}, not the Intel transcendental
implementation. The adapter preserves the older operation names expected by
the project. Successful values are returned unchanged. The square-root
exception repair is narrowly specified: for an ordered interval wholly below
zero, the current \code{Domain_error "sqrt"} is translated to the legacy
\code{Failure "sqrt_I"}. The unchanged project wrapper raises
\code{Unstable}, which the proof predicate handles by subdivision.

\lstinputlisting[firstline=27,lastline=35,firstnumber=27,caption={\code{ocaml/interval_compat.ml}: \code{sqrt_I}}]{../ocaml/interval_compat.ml}

This pathway was required by an observed baseline failure at the argument
\[[-2.854101966249686,-2.854101966249678].\]
It does not replace the argument by
zero, return an empty interval, or accept a box. Other exceptions are not
blanket-translated into successful exclusions. Partial-domain square root is
the backend's existing behavior and is tested separately. The primitive
library's enclosure guarantees remain part of the trusted numerical base;
finite samples do not prove all such guarantees.

The original bytecode stack limit, $16{,}777{,}216$ machine words, is restored
before loading the mathematical source and checked after each job. On the
64-bit baseline it is a maximum of 128 MiB, not a physical reservation. It is
separate from the per-process address-space ceiling. An exhausted width,
stack, time, memory, or output limit produces an incomplete result, never a
proof success.

\section{Triangle area and center geometry}

For squared sides $x_1,x_2,x_3$, define
\[
U(x_1,x_2,x_3)=2(x_1x_2+x_2x_3+x_3x_1)-x_1^2-x_2^2-x_3^2.
\]
The area is $\sqrt{U}/4$. This is \code{ups_I} followed by
\code{area_I}. The center angle opposite $c$ is computed from
$\cos\gamma=(a^2+b^2-c^2)/(2ab)$; conversely the third side comes from
the cosine law. The code uses sharper endpoint evaluations only when it
certifies the signs needed for those monotonicity choices.

\lstinputlisting[firstline=195,lastline=203,firstnumber=195,caption={\code{ocaml/pent.ml}: \code{area_I}}]{../ocaml/pent.ml}

\subsection{Why the lower-endpoint area is a lower bound}
The name \code{areamin_acute} is historical; its intended contract includes
nonobtuse triangles. It evaluates the area expression at the three lower side
endpoints. Those endpoints need not themselves form a nonobtuse triangle.
The following elementary observation explains the conditional contract.

\begin{proposition}
Suppose $X_1,X_2,X_3>0$ are the squared sides of a nondegenerate nonobtuse
triangle. If $0\le z_i\le X_i$, then
\[
U(z_1,z_2,z_3)\le U(X_1,X_2,X_3).
\]
\end{proposition}
\begin{proof}
Maximize $U$ on the closed box $0\le z_i\le X_i$. On a face $z_i=0$,
$U=-(z_j-z_k)^2\le0$, whereas $U(X_1,X_2,X_3)>0$.
Thus a maximizing point of positive value has all coordinates positive.
If two coordinates are below their upper bounds, their partial derivatives
vanish; subtracting the two equations shows that the third coordinate is
zero, a contradiction. If just $z_i$ is below its upper bound, the other two
are $X_j,X_k$, and stationarity gives $z_i=X_j+X_k\ge X_i$ by
nonobtuseness. Hence the maximum is at $(X_1,X_2,X_3)$.
\end{proof}

If the lower-endpoint Heron expression is nonnegative, its square root is a
lower area bound for every intended nonobtuse triangle in the box. If it is
negative, the attempted evaluation must remain unresolved. It is not a
justification to discard the geometric box. This distinction explains why
negative auxiliary radicands may appear even in a successful subdivision
proof.

The suite checks exact right triangles, permutations and scaling, and
numerical nonobtuse fixtures against an independently evaluated Heron formula.
The displayed proposition supplies the mathematical argument; the generated
fixtures check the implementation and cannot replace that argument.

\section{One contact and its inclination angles}

Use the manuscript's directed contact $B\to A$. A vertex of the pointer $B$
lies on an edge of the receptor $A$. The edge coordinate $x_\alpha$ runs from
one receptor endpoint, between $0$ and $2\sigma$; its midpoint is $\sigma$.
The inclination $\theta'$ is associated with the pointer and $\theta$ with
the receptor. They satisfy
\[
\theta+\theta'\equiv\alpha\pmod{2\pi/5}.
\]
These are not ordinary interior angles at an unspecified vertex. Tuple order
and signs must follow the paper's convention throughout.

Let $h=\sigma-x$ and $\psi=\alpha+3\pi/10$. The right triangle from the
receptor center to the contact point has legs $h,\kappa$. The cosine law
therefore gives
\[
\ell(x,\alpha)^2
=1+h^2+\kappa^2-2h\cos\psi+2\kappa\sin\psi.
\]
The production function \code{ell_aux} uses the same construction with an
auxiliary angle rather than this expanded expression. The expansion is useful
as a separate reference in the Wolfram tests. It is not substituted into the
interval calculation, since reassociation would change interval widths.

\lstinputlisting[firstline=247,lastline=258,firstnumber=247,caption={\code{ocaml/pent.ml}: \code{ell_aux}}]{../ocaml/pent.ml}

In particular,
\[
\ell(x,0)^2=4\kappa^2+x^2,
\qquad
\ell(x,\pi/5)^2=(1+\kappa)^2+(\sigma-x)^2.
\]
The first is the familiar slider-contact formula. The second describes the
vertex--edge contact across the shared diagonal of the ice-ray family.

\subsection{Extended contacts and branch boundaries}
The extended domain $0\le\alpha\le4\pi/5$ represents both pointer
directions. For $\alpha>2\pi/5$, exchange the roles using
$(x,\alpha)\mapsto(2\sigma-x,\alpha-2\pi/5)$ and exchange/shift the two
inclinations as in Section~\PaperContactCoordinates. At a box crossing
$\alpha=2\pi/5$, both branches are evaluated and their bounds combined.
A single branch cannot be chosen merely from the midpoint of an interval.

\code{ellxmono_r} includes the stationary candidates as well as the relevant
endpoints. For fixed $x$, the candidate in $\alpha$ is
$\pi/5-\arcsin((x-\sigma)/\sqrt{\kappa^2+(x-\sigma)^2})$.
For fixed $\alpha$, the candidate in $x$ is
$\sigma+\sin(\alpha-\pi/5)$. A candidate whose interval meets the input
box is retained. Extra candidates broaden the result and are harmless;
missing a necessary extremum can invalidate an enclosure.

\section{Pair feasibility and inverse contact domains}

\subsection{The necessary pair test}
\code{Pet.pet} asks whether a box of center distance and inclinations can
contain a nonoverlapping pair. Its Boolean is deliberately asymmetric:
\code{false} is an exclusion, whereas \code{true} only means that the pair
has not been excluded by this test. A wide box may return true even if it
contains no actual nonoverlapping pair.

The function first reduces the inclinations modulo $2\pi/5$, splits at zero,
and exchanges the two roles so that their absolute inclinations are ordered.
It then treats equal-sign and opposite-sign branches. In each branch, the
vertex-to-center distance is compared with the radial boundary of a regular
pentagon. The factors $\kappa/\cos(\cdot)$ are distances to a side in the
appropriate sector. Endpoint or stationary-direction tests bound this radial
comparison. The final result is an existential disjunction over the branch
list: it is enough that one branch remains possible.

\lstinputlisting[firstline=91,lastline=98,firstnumber=91,caption={\code{ocaml/pet.ml}: \code{pet}}]{../ocaml/pet.ml}

The generated tests independently construct the five vertices of each
pentagon and use the separating-axis theorem with an explicit numerical
margin. A comfortably separated pair must not be rejected by
\code{Pet.pet} when put in a small enclosing parameter box. These are
floating-point diagnostics, labeled as such. Contact cases are supplemented
by exact or interval identities; a numerical tolerance is never treated as a
proof of separation at equality.

\subsection{The banana and inverse-banana charts}
The historical names \code{banana} and \code{danana} denote two planar
views of the same directed-contact domain. They are implementation names,
not new assumptions. In the first view the inputs are $(\theta',d)$ and the
output is a range for $\theta$ on the positive-height branch. The negative
branch is obtained by the stated reflection and period shift.

The five boundary curves have the following inverse expressions. Here $u$
is the relevant pointer inclination, and each expression is used only on its
assigned box and branch:
\begin{center}
\begin{tabular}{@{}cl@{}}\toprule
Curve & Center distance as a function of $u$\\\midrule
I & $2\kappa/\cos(u+\pi/5)$\\
II & $2\cos u$\\
III & $2\kappa/\cos(u-\pi/5)$\\
IV & $\cos u+\sqrt{\cos^2u+\kappa^2-1}$\\
V & $(1+\kappa)/\cos u$\\\bottomrule
\end{tabular}
\end{center}
These follow from the inradius right triangle and cosine law at the limiting
contact directions. The domain is divided into five ridgeless pieces and a
sixth piece containing a ridge. \code{find_critical_points} collects corners,
intersections with the rectangle's vertical and horizontal sides, and the
appropriate boundary candidates. The ridge piece has its own extremum rule.
Two global extremal points must be included whenever their intervals meet the
input rectangle.

\lstinputlisting[firstline=150,lastline=171,firstnumber=150,caption={\code{ocaml/banana.ml}: \code{critical_points_banana}}]{../ocaml/banana.ml}

C015 makes the two global-point alternatives independent. In the older
unparenthesized expression, an \code{else} branch could contain the list
concatenation, so one candidate could be masked by the first condition. The
new grouping can only add the missing candidate. A deterministic test chooses
a rectangle meeting both and checks their presence. This is a defensive
correction to the enclosure calculation, not an established counterexample to
an old numerical inequality.

C005 separately changes \code{someinter_I}: only certified interval
disjointness returns \code{None}. An unexpected exception no longer becomes
an empty intersection. The original catch-all was too broad as an interface
contract, even though the baseline jobs happened to complete.

\lstinputlisting[firstline=54,lastline=56,firstnumber=54,caption={\code{ocaml/banana.ml}: \code{someinter_I}}]{../ocaml/banana.ml}

\begin{auditnote}
The trigonometric formulas and their inverse identities are documented and
sample-tested. A complete universal proof of the ridgeless partition and all
critical-point assertions is not claimed here. The remaining trust is stated
explicitly in \code{AUDIT.md}; the added tests include independent witnesses
and boundary cases, rather than only round trips that might hide mutually
compensating mistakes.
\end{auditnote}

\section{Reconstructing a center triangle}

For an oriented center edge $XY$, its complete datum is
\[
(d_{XY},\theta_{XYZ},\theta_{YXZ},\arc_Z).
\]
The first two inclinations belong to the two endpoint pentagons, and
$\arc_Z$ is the opposite center angle. A complete triangle record orders
these edge data as $AB,BC,AC$, preceded by an area lower bound. Do not confuse
that order with the output $(d_{BC},d_{AC},d_{AB})$ of the shared 3C charts.

\subsection{Two contact edges: four variables}
The two contacts at the outer pentagon $B$ are encoded by two extended
$(x,\alpha)$ pairs. The corresponding center distances and inclinations
come from \code{ellthetax}. Their values determine the center angle at $B$
modulo $2\pi/5$:
\[
\arc_B=\frac{2\pi}{5}-(\theta_{BAC}+\theta_{BCA})
       +k\frac{2\pi}{5}.
\]
The possible angle range from the side bounds determines the integer $k$.
\code{getint} returns a unique integer, none, or \code{Unstable} when the
choice is not resolved. It must not choose an arbitrary integer from a wide
box. The cosine law then gives $d_{AC}$, and the remaining angle relations
recover the last two inclinations. A right angle is retained; rejection is
strictly above $\pi/2$.

\lstinputlisting[firstline=298,lastline=325,firstnumber=298,caption={\code{ocaml/banana.ml}: \code{fillout2C}}]{../ocaml/banana.ml}

\subsection{One contact edge: five variables}
The input is
\[
(d_{AB},\theta_{ABC},\theta_{BAC},d_{BC},d_{AC}).
\]
Here $AC$ is required to be a contact edge. First reconstruct the center
triangle and the inclination of $A$ relative to $AC$. Then use both inverse
contact branches to obtain possible inclinations of $C$. The final $BC$ pair
test removes only branches proved impossible. The output is a list of
possible full edge-data records, not a single arbitrary choice.

The wrapper \code{forall_dom} applies a rejection predicate to every branch.
Thus it rejects the parameter box only if \emph{all} represented contact
branches are rejected. The key functions instead take the union of keys of
all branches. These quantifiers are deliberately different.

\subsection{No required contact: six variables}
The input is
\[
(d_{AB},\theta_{ABC},\theta_{BAC},d_{BC},\theta_{CBA},d_{AC}).
\]
The three sides determine the center triangle. At each center,
\[
\arc_A+\theta_{ABC}+\theta_{ACB}\equiv0\pmod{2\pi/5},
\]
and similarly at $B,C$. These equations determine the remaining inclinations.
Three pair tests supply necessary nonoverlap constraints. No contact is
imposed, so this chart includes the contact subspaces and is appropriate for
the general central triangle in the 6D jobs.

The helper's area value is conditional on the intended nonobtuse subset.
Additional irrelevant obtuse inputs in a box do not require their own area
bound; what matters is that no intended nonobtuse witness is lost. A helper
returning \code{None}, however, must justify that the relevant represented
set is empty.

\section{Two triangles, their contact types, and squeezing}

\subsection{The shared-edge relation}
For adjacent triples $(A,B,C)$ and $(A,D,C)$, the common pentagons $A,C$
are the same objects. In the shared 3C charts, their contact relation gives
\[
\bar\beta_D=2\pi/5-\bar\beta_B,\qquad
\bar x_{\beta,D}=2\sigma-\bar x_{\beta,B}.
\]
The remaining chart parameters are independent until the geometric tests
restrict them. The code evaluates the first triangle, uses its subcritical
and shared-longest assumptions, then evaluates the second with these exact
shared-variable identities.

The eight names, in their fixed order, are
\code{pint,pinw,lj1,lj2,lj3,tj1,tj2,tj3}. Different names may describe different
choices of shared edge within the same geometric contact type. They must not
be merged solely because the pictures look alike.

\subsection{What the dimer groups actually exclude}
Group 1 treats the two isosceles endpoints and four coordinate restrictions.
Group 2 contains 64 calls covering 48 distinct ordered pairs; the original
repeated calls are retained. These use no near-ice-ray alternative. Group 3
contains the other sixteen ordered pairs. Twelve use \code{pseudo=false};
four additionally allow the explicitly coded pseudo-dimer alternative.
Every group-3 case retains the $0.01$ neighborhood alternative.

Group 4 returns to the four exceptional ordered pairs with
\code{pseudo=false}, no neighborhood alternative, and area cutoff
$2\ak-0.007$. Since
\[
2\ak-0.007>2\ak-0.008,
\]
this is a stronger lower bound than needed on these terminal cases. The
constant is not changed merely to agree typographically with $\eM$.
It does not by itself cover the separate isosceles-peripheral cases.

Group 5 historically bundled squeezing and four \code{one127} cases.
The release exposes squeezing as its own proof job and \code{one127} as an
optional historical diagnostic. The latter's existing domain and conclusion
are retained, but its relevance to the current proof is not asserted.

\subsection{Four contact-sign calculations}
The proof of Lemma~6.5 (\code{lemma:squeeze-area}) uses the area cap
$\ak+\eN$, not $\amin+\eN=\ak$. It first needs the signs of the supporting
contact lines, then a separate inequality to prove area decrease. The new
module \code{squeeze_sign.ml}, navigation code \code{GJULAOH}, establishes
the former on the complete four-dimensional domain.

Fix the center of $B$ at the origin and one of its vertices at $(1,0)$.
There are four pointer choices: $B\to A,B\to C$; $B\to A,C\to B$;
$A\to B,B\to C$; $A\to B,C\to B$. In each case use the two contact
coordinates $(x_\alpha,\alpha)$ and $(x_\gamma,\gamma)$ of Appendix A.1,
with $0\le x_\alpha,x_\gamma\le2\sigma$ and
$0\le\alpha,\gamma\le2\pi/5$. Closed endpoints retain sliders and
one-sided vertex--vertex contact representations. A relative rotation
$2\pi m/5$, $m=0,\ldots,4$, determines the second attachment. Each of the
five alternatives is checked. Although at a fixed feasible point the acute
angle interval selects at most one alternative, an interval box need not
have a unique index. No midpoint-based branch selection is used.

Write $\mathbf e(\theta)=(\cos\theta,\sin\theta)$ and let $J$ be a positive
quarter turn. For $B$ pointing to a neighbor $X$, put
$h=x-\sigma$, $\phi=\alpha-\pi/5$. The neighbor center, normal from $B$
toward $X$, and neighbor vertex orientation are
\[
\mathbf n=\mathbf e(\phi),\quad
\mathbf c_X=(1,0)+\kappa\mathbf n+hJ\mathbf n,\quad
\omega_X=\alpha+3\pi/5.
\]
For $X$ pointing to $B$, put $\psi=2\pi/5-\alpha$. Exchanging the pointer
and receptor in the same construction gives
\[
\mathbf n=\mathbf e(\pi/5),\quad
\mathbf c_X=\mathbf e(\psi)+\kappa\mathbf n+hJ\mathbf n,\quad
\omega_X=\pi+\psi.
\]
Both formulas use the endpoint definition of $x$ in the paper, not a
midpoint offset. Rotating the entire $C$ record through $2\pi m/5$
completes the chart. The unit four-cube used in subdivision is mapped to
these ranges by interval scaling, without truncating a parameter boundary.

Let $\mathbf a=\mathbf c_A$, $\mathbf c=\mathbf c_C$ and
$\mathbf g=\mathbf c-\mathbf a$. The chosen orientation has
$H=\det(\mathbf a,\mathbf c)>0$, so $B$ lies above $AC$. Define
\[
q_A=\mathbf n_A\cdot\mathbf g,\quad
r_A=\det(\mathbf n_A,\mathbf g),\qquad
q_C=\mathbf n_C\cdot\mathbf g,\quad
r_C=\det(\mathbf n_C,\mathbf g).
\]
In the orthonormal frame along $AC$, the tangent $J\mathbf n$ has
components $(r,q)/\lVert\mathbf g\rVert$, and the inward normal toward
$B$ has upward component $r/\lVert\mathbf g\rVert$. Hence
\[
-q_A>0,\quad r_A>0,\quad q_C>0,\quad r_C>0
\]
prove the required negative/positive finite slopes and the supporting-side
assertions together. They also exclude horizontal and vertical normals.
No division by an uncertain slope is involved.

The classifier first removes the wrong orientation, then tries, in order:
insufficient $AC$ center separation; a certified obtuse angle; area above
$\ak+\eN$; the four desired signs; and strict interior overlap of $A,C$.
The last test uses actual vertices and all ten edge-normal axes. The
projection of a unit regular pentagon on its own outward edge normal is
$[-1,\kappa]$. Directed endpoint comparisons certify strict overlap of the
projection interiors; all axes must pass. Touching cases are not discarded.
Any unresolved box is bisected. A numerical exception or a subdivision limit
never becomes an accepted box.

The binary-forest checks require all five roots to be visited and
$\text{nodes}=2\,\text{splits}+5$,
$\text{leaves}=\text{splits}+5$. Each of the four pointer jobs has a separate
process and log. The detailed chart derivation and separating-axis argument
are also in \code{docs/SQUEEZING.md}.

\subsection{The squeezing area inequality}
Let $\ell=d_{AC}$ and $h=d_{AB}\sin\arc_A$. The signed-normal calculation
just described supplies positive finite quantities
$\sigma_A=-q_A/r_A$ and $\sigma_C=q_C/r_C$. With $A$ fixed and
$\ell'=-1$, the contact equations give
\[
x_B'=-\frac{\sigma_C}{\sigma_A+\sigma_C},\qquad
h'=\frac{\sigma_A\sigma_C}{\sigma_A+\sigma_C}>0.
\]
A horizontal change of frame gives exactly the motion directions used in the
paper. The area derivative is $(-h+\ell h')/2$, so decrease is equivalent to
\[
(1/\sigma_A+1/\sigma_C)h>\ell.
\]
The existing \code{squeeze_calc} accepts this condition or either of
$\sigma_A<h/\ell$, $\sigma_C<h/\ell$. For positive finite slopes, either
of the latter implies the reciprocal inequality. The disjunction is thus a
sufficient interval test on every geometric branch retained by the sign
calculation. The sign searches and this inequality calculation are both
selected by \code{--claim lemma:squeeze-area} (\code{LZMJIZU}).

\lstinputlisting[firstline=270,lastline=290,firstnumber=270,caption={\code{ocaml/dimer.ml}: \code{squeezable}}]{../ocaml/dimer.ml}

\section{The exact ice-ray shear identity}

The final convexity calculation is no longer needed as a theorem premise.
The manuscript now has a shorter geometric argument. Fix $A$ and all
orientations. Let $\mathbf e,\mathbf f$ be the orthonormal directions along
and perpendicular to the shared diagonal at the ice-ray position. The shear
parameter is the signed offset from the receptor-edge midpoint, not the
absolute contact parameter. Thus $x=\sigma+t$.

The slider directions give
\[
\begin{split}
\cc_C(t)-\cc_C(0)&=t\mathbf f,\\
\cc_B(t)-\cc_B(0)&=t(-\sigma\mathbf e+\kappa\mathbf f),\\
\cc_D(t)-\cc_D(0)&=t(\sigma\mathbf e+\kappa\mathbf f).
\end{split}
\]
The equal signed slide lengths follow from the small displacement triangle
with angles $\pi/5,2\pi/5,2\pi/5$. Therefore the diagonal vectors satisfy
\[
\mathbf u(t)=\mathbf u(0)+t\mathbf f,\qquad
\mathbf v(t)=\mathbf v(0)-2\sigma t\mathbf e.
\]
At $t=0$ they are parallel to $\mathbf e,\mathbf f$, respectively. Taking
half their determinant yields
\[
\boxed{\area\{T_0(t),T_1(t)\}=2\ak+\sigma t^2.}
\]
The minimum is unique at zero. No periodic extension is needed for this local
area computation.

For an independent algebraic check, the squared sides of one triangle are
\[
4\kappa^2+x^2,\qquad
4\kappa^2+\frac{x^2}{4\kappa^2},\qquad
(1+\kappa)^2+(\sigma-x)^2.
\]
Their Heron polynomial is $4B(x)^2$, where
\[
B(x)=\sigma x^2+(\kappa-1)x+\sigma(4\kappa+1).
\]
Its leading coefficient is positive and its discriminant is
$-(13+15\sqrt5)/8<0$, so the area is $B(x)/2$. Adding the partner with
parameter $2\sigma-x$ gives the same identity. Both checks are included in
the Wolfram suite; the exact polynomial identities were also checked with
SymPy during preparation.

The optional \code{one_dimer2} diagnostic uses the absolute parameter
$x$ on $\sigma+[-0.1,0.1]$, corresponding to the centered offset
$|t|\le0.1$. Its greater width is conservative. This consistency check is not
a premise of the geometric area proof.

\section{Automatic differentiation along the other paths}

A first-order jet is $(f,f')$; a second-order jet is $(f,f',f'')$, not a
factorial-scaled coefficient tuple. The operations are the chain and product
rules applied with interval arithmetic. For example,
\[
(fg)''=f''g+2f'g'+fg'',\qquad
(f^{-1})''=2(f')^2/f^3-f''/f^2.
\]
The implementation does not estimate derivatives by finite differences.
Those estimates, when used in tests, are only diagnostics.

\lstinputlisting[firstline=424,lastline=428,firstnumber=424,caption={\code{ocaml/autodiff.ml}: \code{sqrtD2}}]{../ocaml/autodiff.ml}

The two twist calculations handle the LJ charts before the final common
shared-edge path. The path checks then cover a $5\times5$ ordered chart
matrix. Their original first-derivative seeds and initial boxes are retained;
the new driver names each pair separately instead of launching all pairs as a
side effect of loading \code{autodiff.ml}. This makes both ordinary testing
and selective rerunning possible.

The thirty-one original numerical test intentions are now retained as named
thunks, including the previously shadowed first- and second-order product
tests. The suite also checks exact jets for a cube, reciprocal, square root,
sine and cosine, and checks the exact shear identity at several values of
$t$. Equalities involving transcendental interval results are tested by
necessary enclosure/overlap relations, not by imposing an arbitrary equality
of rounded endpoints.

\begin{auditnote}
The code's derivative values are only useful on the correct parameter path.
The appendix's chart order and the shared-angle identities must remain the
basis of that interpretation. The present commentary specifies these links
and the suite exercises them, but does not claim a fresh universal proof of
every first-derivative chart enclosure.
\end{auditnote}

\section{Meet in the middle: why a feasible cluster must survive}

Let a counterexample consist of a central triangle and one or more peripheral
triangles. Along a shared edge, the common data are the length and two
inclinations. Rather than store every central/peripheral combination, the
algorithm stores peripheral boxes in a hash table keyed by discretizations
of those three shared variables. Central boxes retrieve lower area bounds
from matching keys.

\subsection{Closed keys and orientation}
For width $w>0$ and closed interval $I$, the scalar key list contains every
integer from $\lfloor\inf(I/w)\rfloor$ through
$\operatorname{ceil}_{\rm strict}(\sup(I/w))-1$. The strict ceiling
returns $n+1$ when its input is the integer $n$. Hence a singleton at a cell
boundary shares a key with either neighboring interval containing that point.
For triples, the Cartesian products are hashed. Hash collisions can create
extra matches, but cannot destroy a necessary match.

\lstinputlisting[firstline=30,lastline=36,firstnumber=30,caption={\code{ocaml/meet-in-middle.ml}: \code{make_a_key}}]{../ocaml/meet-in-middle.ml}

Adjacent triangles traverse the common edge oppositely, so the central
key routine negates both inclinations before reducing them modulo $2\pi/5$.
Periodization keeps all endpoint representatives. Key tests cover negative
intervals, exact grid points, wrap boundaries, and generated dyadic intervals
with a known common point.

\subsection{Lower area bounds and upper cuts}
For every key, the stored peripheral area is a \emph{minimum lower bound}
over compatible boxes. A hash collision can lower this bound; that makes
exclusion more difficult, not unsound. If no compatible entry remains, the
central box can be excluded, provided the key coverage property holds.

Suppose a central lower bound is $a_0$, and compatible peripheral lower
bounds are $a_1,\ldots,a_m$. If
\[
a_0+a_1+\cdots+a_m>A,
\]
where $A$ is the counterexample's total-area cap, that central box is outside
the domain. Otherwise a peripheral participating in this central possibility
must obey
\[
\area(T_i)\le A-a_0-\sum_{j\ne i}a_j.
\]
The implementation stores an outward upper endpoint of this expression as an
allowable cut. When several central possibilities use one peripheral key,
it takes the \emph{maximum} allowable cut. A minimum would be incorrect:
one feasible central alternative could be removed by another incompatible
one's tighter requirement.

\lstinputlisting[firstline=87,lastline=93,firstnumber=87,caption={\code{ocaml/meet-in-middle.ml}: \code{recut}}]{../ocaml/meet-in-middle.ml}

\begin{contract}
Assume each geometric fill function encloses every intended witness or proves
emptiness; each outside-domain predicate is one-sided sound; each area value
is a lower bound; and matching witness data share a key. Then every feasible
cluster of total area at most $A$ survives each pruning round.
\end{contract}
\begin{proof}
Its peripheral boxes cannot be removed initially by their geometric tests.
Their keys give central matches. Each stored peripheral minimum is no greater
than the corresponding witness area, and the central lower bound is no greater
than its witness area. Thus their sum cannot exceed $A$ strictly. The cut
propagated by this central witness is at least the area of its peripheral
witness. Taking a maximum over possible central matches cannot lower that
cut. The witness therefore survives the next peripheral pruning. Subdivision
covers the closed parent boxes, so the same argument applies inductively.
\end{proof}

This proof separates the combinatorial pruning invariant from the geometric
fill contracts. It does not assert that a flawed geometric \code{None} can be
repaired by the hash-table logic.

\subsection{Termination and tests with feasible examples}
The key width halves at each outer step. A successful calculation ends only
when the relevant peripheral set or the central set is empty. Reaching the
minimum width raises an exception. There is no rule ``small enough means
verified.'' Independent workers have separate tables and result files.

The tests do not use only impossible examples. They construct small finite
cluster problems, enumerate all compatible combinations directly, and
compare that independent enumeration with \code{mk_one_cencoord}. A feasible
combination exactly at the target must survive, with both peripheral cuts at
least its actual areas. Impossible fixtures must be removed. Additional tests
force unresolved fills and width exhaustion; neither is allowed to become a
pass. Duplicate-key and collision fixtures check conservative minima and
maximum cuts.

\section{The eight principal computational inputs}

In this section, the stated geometric reductions belong to the paper.
The commentary identifies their terminal computational domains; it does not
replace those reductions by the mere existence of source functions. Every
area rejection is strict. Closed equality faces are retained unless a
separate geometric argument excludes them.

\subsection{Shared edge below 1.8: \code{NKQNXUN}}
Lemma~\PaperPseudoOne negates a shared-edge upper bound. The total target is
$2\ak$, and the central cap is $\ak+\eN$. There are three terminal
branches: outside double contact, a midpoint endpoint, and a slider endpoint.
The code uses the retained nonshared-length consequence rather than the
combinatorial arrow itself, so the added tied-arrow relaxation is retained
at this interface. The closed ranges include the boundary 1.8.

The functions are \code{precluster1_case0}, \code{precluster1_case1}, and
\code{precluster1_case2}. The first uses the four-dimensional two-contact
chart; the latter two use a five-dimensional contact-central chart. The
side-range cuts are supported by explicit interval assertions in the source
and independent exact-number checks in the Wolfram suite.

\subsection{A nonshared longest edge above 1.8: \code{RWWHLQT}}
For Lemma~\PaperPseudoTwo, fix the shared pentagons and the peripheral.
The peripheral data remain O2C, subcritical, and shared-longest; dropping the
selected arrow does not drop that last length consequence.
In the paper labels, let $p=d_{AB}$ be the primary-contact side,
$q=d_{BC}$ the other nonshared side, and $s=d_{AC}$ the shared side.
The compact minimum retains
\[
p,q\le1.8,\qquad \max(p,q)\ge s.
\]
Besides O2C, a minimum has midpoint or slider contact and lies on a tie or a
nonshared length-1.8 boundary. The primary-contact descent fixes the larger
incident distance, so it preserves the lower bound on the maximum.

In the two five-dimensional code branches the labels are permuted:
\[
(AB,AC,BC)_{\rm code}=(AC,AB,BC)_{\rm paper}.
\]
Thus code $AB$ is shared and code $AC$ is contact. The boundary helper deliberately
retains all three equalities and both nonshared 1.8 boundaries, a superset of
the required stopping events:

\lstinputlisting[firstline=451,lastline=457,firstnumber=451,caption={\code{ocaml/mitm-calcs.ml}: \code{pseudo2_boundary_excluded}}]{../ocaml/mitm-calcs.ml}

Both \code{precluster2_case1} and \code{precluster2_case2} call this actual
helper; the tests exercise the same helper, not a duplicate expression.
The midpoint/slider filters and all other bounds remain in their respective
functions. The O2C branch \code{precluster2_case0} uses the separate two-contact chart. All terminal branches belong to the required proof set.

\subsection{Pseudo-dimer pair area: \code{BXZBPJW}}
The target in Lemma~\PaperPseudoArea is $2\ak-\eM$. The two mixed-contact
branches use \code{precluster4_case0_experimental} and
\code{precluster4_case1_experimental}, distinguishing which side of the
peripheral equals the shared side. Their central cap is
$\ak+\eN-\eM$. The word \code{experimental} is retained as part of the
historical stable name; it is not evidence of obsolescence.

The triple-contact route uses the earlier type exclusions and the four
stronger $0.007$ cases in \code{Dimer.run_group4}. No attempt is made to
weaken that stronger bound to match the paper. All these comparisons are strict, with interval endpoints directed outward.

\subsection{A pseudo-dimer and its outgoing peripheral: \code{JQMRXTH}}
Lemma~\PaperPseudoThree has target $3\ak+\eM$. The exceptional outgoing
peripheral is not assumed subcritical. In a counterexample,
\[
\area(T_-)<3\ak+\eM-(2\ak-\eM)=\ak+2\eM.
\]
The common initializer uses that cap, which follows from the stated counterexample assumptions.

\lstinputlisting[firstline=585,lastline=593,firstnumber=585,caption={\code{ocaml/mitm-calcs.ml}: \code{init2Cps_precluster6}}]{../ocaml/mitm-calcs.ml}

The shared initializer serves the ordinary incoming and exceptional outgoing
roles. Its first disjunct retains the ordinary constraints where the shared
edge is strictly below 1.8; the second cap is common. This reuse may admit
extra peripheral candidates. Extra candidates make the exclusion harder and
are not a defect. The six jobs distinguish the six ordered choices of
incoming and outgoing edges:
\[
(AB,BC),(AB,AC),(BC,AB),(BC,AC),(AC,AB),(AC,BC).
\]
Their central cap remains $\ak+\eN$ and their total target remains unchanged.
All six cases are selected by the lemma identifier.

\subsection{The isosceles dimer endpoint: \code{KUGAKIK}}
\code{dimer_isosceles_precluster} represents the long-isosceles O2C
peripheral and the central double-contact chart, with total cap $2\ak$.
The peripheral initializer imposes the appropriate shared-side equality and
retains the third-side longest comparison. The main lemma is
Lemma~\PaperIsosceles. This is separate from the local isosceles exclusions
in dimer group 1; both associations remain visible in the index.

\subsection{Three nonobtuse triangles: \code{HUQEJAT}}
Lemma~\PaperThree has two subcritical peripherals attached to distinct edges
of a nonobtuse central triangle. The total cap is $3\ak+\eM$; subtracting
two $\amin$ bounds gives central cap $\ak+2\eN+\eM$.
When the center acquires a contact, the five-dimensional cases distinguish
the edge without a peripheral:
\begin{center}\small
\begin{tabular}{@{}lll@{}}\toprule
Function suffix & Nonshared edge & Peripheral keys\\\midrule
\code{AB} & $AB$ & $BC,AC$\\
\code{BC} & $BC$ & $AB,AC$\\
\code{AC} & $AC$ & $AB,BC$\\\bottomrule
\end{tabular}
\end{center}
The remaining terminal alternative has a general six-dimensional central
chart and two isosceles peripheral charts. It is
\code{calc_pent3_6D_postcluster}. Its historical run has completed; it is
not an unresolved execution issue. The three contact-based alternatives do
not replace it. The nonshared upper side bound 2.02 is checked against the
central cap by an explicit source assertion.

\subsection{Four nonobtuse triangles: \code{QPJDYDB}}
\code{calc_pent4_6D_postcluster} uses a general central triangle, one
subcritical O2C peripheral on each edge, target $4\ak$, and central cap
$\ak+3\eN$. The central side ordering is a symmetry reduction with
non-strict boundaries retained. The six-dimensional chart includes central
contact subspaces, so the explicitly superseded five-dimensional alternative
is removed from the active release rather than kept as a required second
proof.

\subsection{A non-large angle: \code{FHBGHHY}}
The five existing \code{calc_pent2_postcluster_case0} through \code{case4}
cover O2C or one contact together with the distinguished nonshared side at
length 1.8, according to the normalized contact/shared-edge placement.
The target is $2\ak+\eM$ and the central cap is $\ak+\eN+\eM$.
The paper's geometric argument excludes right triangles: the shared
edge is a longest side of the subcritical peripheral, hence greater than
1.72 and less than $2.1<\kappa\sqrt8$. It is therefore a leg, not the
hypotenuse, of any right central triangle. Its other leg is at least
$2\kappa$, giving
\[
\area(T_0)+\area(T_1)>\amin+1.72\kappa>2\ak+\eM.
\]
No new right-angle job is required for that reason. The new manuscript note
asks only that the leg be identified and the compact minimum be stated
instead of invoking termination of a possibly infinite translation path.

\section{Tests, execution evidence, and the trust boundary}

The OCaml suite exercises the actual release functions. It includes exact
dyadic examples, strict-boundary and exception tests, contact-chart fixtures,
first/second derivative rules, and finite MITM problems with an independent
exhaustive answer. Generated diagnostics use recorded seeds and independent
polygon construction. Old work counts, rounded endpoints, and historical
runtimes are not mathematical acceptance criteria. The 153 ordinary contract
checks, the extended generated groups, and the 71 contact-sign checks have
separate stage results.

The Wolfram package defines its functions without executing them on import.
Its 161 named checks are documented individually in
\code{docs/WOLFRAM_CHECKS.md}: exact assertions from the paper, supporting
identities, implementation fixtures, and numerical diagnostics have different
roles. The report includes the displayed paper number, stable label,
mathematical assertion, and reason for the check. Numerical examples do not
replace uniform interval searches.

Python tests exercise scheduling, memory admission, failure handling,
checkpoint identities, report copying and confirmation with synthetic child
processes. They are not executions of OCaml or Wolfram mathematics.

The interval backend is trusted for outward rounding, including
whole-interval transcendental ranges. The compiler/runtime and application
coordinate routines are also trusted components. The formulas and retention
arguments in this commentary, independent geometric fixtures, finite
exhaustive tests, and exact symbolic checks provide complementary scrutiny;
this is not a machine-checked formalization of every implementation branch.

\subsection{Source-specific confirmation}
The release requires 70 proof jobs and three function-test stages. Four
additional historical or consistency diagnostics are optional. The
controller requires normal exit, an exact job marker, the retained stack
setting, and unchanged source identities. A failure, timeout, interruption,
missing job or missing marker cannot become a pass. An intact checkpoint is
reused only for identical source and installed-environment fingerprints.

The command \code{python3 confirm.py --check} validates the complete required
OCaml set and the full Wolfram export, then writes
\code{verification/CONFIRMATION.md}. The record is an execution summary, not
a formal proof certificate or a signature identifying an operator. A
source-matching verified distribution can be created with
\code{confirm.py --bundle}. Historical evidence is recorded separately and
cannot fill missing jobs in the current confirmation.

\subsection{Resources and long-term reproducibility}
A recorded 2026 Intel/WSL execution of the general six-dimensional
three-triangle calculation took about 6 hours 36 minutes wall time. The
largest observed individual resident peak in that concurrent baseline was
about 2.58 GiB. These are measurements, not runtime predictions. Four
independent workers have separate hash tables; the controller budgets
4 GiB address space per worker and 8 GiB additional memory headroom. It does
not treat swap as available admission memory.

Preserve the exact interval source, its notices, compiler/package metadata,
and run records. A package name, development pin, or container image is not
by itself an archive of all required source. The environment guide gives a
version-pinned installation and an explicit source-archiving command.
The mathematical source, parameter domains, and evidence remain inspectable
without a private directory layout or an unpublished notebook.

\end{document}
